\section*{第 \ref{ch:g8:circuit-laws} 章 --- 电路中的电流与电压定律}

\begin{solution}{exo:g8:circuit-laws:1}
串联电流定律：一条回路上，处处电流强度相同。
节点定律：流入节点的电流之和等于流出的电流之和。串联电压
定律：电池电压等于沿回路各用电器分得的电压之和。
并联电压定律：跨接于相同两节点的各支路，
电压相同。
\end{solution}

\begin{solution}{exo:g8:circuit-laws:2}
$0.8 - 0.55 = \qty{0.25}{A}$.
\end{solution}

\begin{solution}{exo:g8:circuit-laws:3}
$9 - 2.5 - 3.8 = \qty{2.7}{V}$.
\end{solution}

\begin{solution}{exo:g8:circuit-laws:4}
各为 \qty{12}{V}，即并联电压定律。同一定律
也使家庭电路中每个插座的电压都相同，均为 \qty{230}{V}，
即并联线路各处电压相等。
\end{solution}

\begin{solution}{exo:g8:circuit-laws:5}
以 $I_C$ 为例：已知 $I_{\text{电池}}$ 和 $I_B$ 后，
节点定律即可确定 $I_C = \qty{0.2}{A}$——接在此处的任何电流表都只能证实这一结果，
（同理，$U_C$ 由并联电压定律在
已知 $U_B$ 后确定）。
\end{solution}

\begin{solution}{exo:g8:circuit-laws:6}
\qty{2.7}{V} 的灯泡：电流相等时，电压分配越大，
说明它对电流的阻碍作用越强。下一章将测量
这种阻碍作用，即电阻。
\end{solution}

\begin{solution}{exo:g8:circuit-laws:7}
$2.2 + 1.9 + 1.4 = \qty{5.5}{V}$，而电池电压为 \qty{6}{V}：
相差半伏，超出此类读数的正常误差范围；
必有一个表盘读错了，或有一段电压未被测到。
\end{solution}

\begin{solution}{exo:g8:circuit-laws:8}
$0.4 + 0.3 + 0.3 = \qty{1.0}{A}$ 流出，而流入仅 \qty{0.9}{A}，
分叉处“凭空产生”了 \qty{0.1}{A}——
节点定律即可判定其有误。
\end{solution}

\begin{solution}{exo:g8:circuit-laws:9}
回路中各用电器的电压之和必须等于电池的 \qty{4.5}{V}：
断开的开关 \qty{4.5}{V} $+$ 不亮的灯泡 \qty{0}{V} $= 4.5$——
平衡。（开关闭合后，电压重新分配：开关两端电压接近零，
灯泡两端电压接近 \qty{4.5}{V}，仍然平衡。）
\end{solution}

\begin{solution}{exo:g8:circuit-laws:10}
变化前：$3 \times 1.5 = \qty{4.5}{V}$。变化后：$4.5 - 1.5 =
\qty{3.0}{V}$——一节电池反接，电压相减。灯泡
由 \qty{4.5}{V} 时的亮度变为 \qty{3}{V} 时的较暗亮度。
\end{solution}

\begin{solution}{exo:g8:circuit-laws:11}
支路一跨接电池：$U_E + U_F = \qty{6}{V}$（沿支路的串联电压定律），
故 $U_F = 6 - 2.2 =
\qty{3.8}{V}$。支路二跨接相同两节点：$U_G =
\qty{6}{V}$（并联电压定律）。
\end{solution}

\begin{solution}{exo:g8:circuit-laws:12}
$I_R = 0.6 - 0.45 = \qty{0.15}{A}$（节点定律）。$U_R = U_Q
= \qty{2.3}{V}$（并联电压定律）。$U_P = 6 - 2.3 =
\qty{3.7}{V}$（串联电压定律）。一个可被判错的说法：
“$I_R = \qty{0.2}{A}$”——节点处的电流对不上：
$0.45 + 0.2 = 0.65 \neq 0.6$；或者“$U_R = \qty{2.6}{V}$”——
并联支路的电压不能互相矛盾。
\end{solution}

\begin{solution}{pb:g8:circuit-laws:1}
\textbf{1.} 串联。串联时两只灯泡依次阻碍电流，
回路中的电流远小于单只灯泡的
\qty{0.4}{A}——测得的 \qty{0.2}{A} 与此相符。
（并联则是两条完整支路的电流\emph{相加}。）
\textbf{2.} 并联：两条支路各约 \qty{0.4}{A}，
电池输出约 \qty{0.8}{A}。
\textbf{3.} 串联电压定律将 \qty{6}{V} 分到
两只相同灯泡上，每只约 \qty{3}{V}。
\textbf{4.} 并联——每只灯泡都将跨接整个
\qty{6}{V}。从观察窗看，串联的一对灯泡发光较暗：
每只灯泡两端的电压只有平常的一半。
\textbf{5.} $I_Z = 0.5 - 0.35 = \qty{0.15}{A}$（节点定律）；
$I_X = \qty{0.5}{A}$——X 位于未分叉的通路上（串联
电流定律）。
\textbf{6.} $U_{\text{一对}} = 6 - 2.8 = \qty{3.2}{V}$（回路电压
定律）；$U_Y = U_Z = \qty{3.2}{V}$（并联电压定律）。
\textbf{7.} 电压相同，电流路径不同：两只灯泡对电流的
阻碍并不相同——在相同电压下通过电流较小的 Z，
阻碍作用更强（路更难走）。
\textbf{8.} 变化：去掉 Y 所在支路即少了一条通路，
剩下的 Z 通路对电流的阻碍比原来两者共同的阻碍还要大，因此
电池输出的总电流变小。
\textbf{9.} 第一条：回路定律，$2.5 + 4.1 =
\qty{6.6}{V}$，而电源电压只有 \qty{6}{V}。第二条：
节点定律：流入 \qty{0.3}{A}，流出 $0.2 + 0.15 =
\qty{0.35}{A}$。
\textbf{10.} $6 - 2.5 = \qty{3.5}{V}$。
\textbf{11.} $0.2 + 0.15 = \qty{0.35}{A}$。
\textbf{12.} 因为四条定律对任何真实电路都相互约束：
每个读数都须与其他读数的和相符，因此一个
编造的数字至少会违反一条定律，
仅凭算术即可判定其有误。
\textbf{13.} 例如：“同一条路，同一股电流——串联电流定律。
在每个分叉处，流进与流出平衡——节点定律。
沿任一回路，电池的“水流”
被各用电器耗尽，一滴不剩——串联电压定律。
挂在同样两个节点之间的各支路，
都从同一高度取水——并联电压定律。”
\end{solution}
